\chapter{Analyse critique, apports et perspectives}

\section{Qualite logicielle}

La qualite logicielle de SupportFlow se lit a plusieurs niveaux. Au niveau du code, le decoupage en couches et en services specialises constitue deja un socle solide. Au niveau fonctionnel, la qualite se mesure a la coherence des parcours. Au niveau systeme, elle depend de la capacite a relancer l'environnement complet et a retrouver un etat credible de demonstration.

\section{Maintenabilite}

Le projet est relativement maintenable grace a :

\begin{itemize}[leftmargin=1.2cm]
  \item une architecture modulaire ;
  \item une documentation abondante ;
  \item un environnement Docker reproductible ;
  \item des scripts de test et de diagnostic ;
  \item des conventions metier de plus en plus explicites.
\end{itemize}

La maintenance doit cependant rester vigilante sur les points de jonction entre les composants. Une regression sur un champ de tri, sur une permission ou sur un identifiant de ticket peut provoquer des erreurs visibles importantes meme si le coeur technique du systeme reste sain.

\section{Analyse des risques techniques}

\begin{longtable}{>{\raggedright\arraybackslash}p{4.4cm} >{\raggedright\arraybackslash}p{3cm} >{\raggedright\arraybackslash}p{5.1cm}}
\toprule
Risque & Impact & Reponse possible \\
\midrule
Lenteur de l'IA locale & Mauvaise experience utilisateur & Timeout court, fallback et asynchronisme. \\
Desynchronisation base / Camunda & Incoherence de supervision & Bootstrap Camunda et verification Cockpit. \\
Permissions incoherentes & Refus ou exces d'acces & Centralisation des autorisations. \\
Seed trop pauvre ou errone & Demo peu representative & Jeu de donnees realiste et controlee. \\
Erreur de mapping frontend / backend & Erreurs serveur evitables & Alignement strict des contrats et champs. \\
Traitements lourds dans le chemin critique & Latence ressentie & Deporter les taches lentes. \\
Regression multi-role & Parcours casses pour certains profils & Tests role par role. \\
\bottomrule
\end{longtable}

\section{Analyse des risques metiers}

Les risques metiers sont parfois plus graves qu'un bug purement technique. Un manager qui ne comprend pas pourquoi un agent n'apparait pas dans une liste, un client qui ne comprend pas pourquoi son ticket reste bloque, ou un agent qui voit une action qu'il ne peut pas executer sont autant de situations qui diminuent la confiance dans l'outil.

\section{Observabilite et supervision}

L'observabilite est bien servie par le projet :

\begin{itemize}[leftmargin=1.2cm]
  \item Camunda Cockpit permet de verifier les instances et les taches ;
  \item les healthchecks Docker indiquent l'etat des composants ;
  \item les logs aident au diagnostic ;
  \item les scripts de test facilitent la verification des parcours.
\end{itemize}

\section{Qualite de demonstration}

Dans un contexte de soutenance, la qualite de demonstration est un critere de succes. SupportFlow repond bien a ce besoin car il ne demande pas seulement d'expliquer des concepts ; il permet de les montrer. On peut ouvrir le frontend, observer Keycloak, verifier les workflows dans Cockpit, consulter la GED, voir les notifications et relancer l'ensemble via Docker.

\section{Maintenance evolutive}

La maintenance evolutive pourrait s'orienter vers :

\begin{itemize}[leftmargin=1.2cm]
  \item une supervision manager encore plus orientee files d'attente ;
  \item des regles metier plus explicites sur les tickets en attente ;
  \item une base de connaissances plus structuree ;
  \item des tests end-to-end plus nombreux ;
  \item une observabilite plus approfondie des traitements asynchrones.
\end{itemize}

\section{Apports techniques et professionnels}

Au niveau technique, le projet a permis de consolider plusieurs acquis majeurs : conception d'APIs REST, decoupage frontend/backend, modelisation des roles et permissions, orchestration BPMN, utilisation d'une GED, deploiement conteneurise, integration continue et demarche GitOps. L'etudiant a egalement acquis une meilleure maitrise des strategies de diagnostic, notamment lorsqu'un probleme ne vient pas d'un composant unique mais de leur interaction.

Au niveau professionnel, ce stage a renforce la capacite a raisonner en termes de valeur metier, de priorisation, de soutenabilite et de preuve de fonctionnement. Il a aussi developpe des reflexes importants : documenter, verifier, rejouer les scenarios critiques, preparer des jeux de donnees coherents et ne pas confondre livraison de fonctionnalites et stabilisation globale du produit.

\section{Analyse critique et perspectives}

\section{Forces du projet}

SupportFlow presente plusieurs forces evidentes :

\begin{itemize}[leftmargin=1.2cm]
  \item une architecture complete et credibilisante ;
  \item une bonne profondeur fonctionnelle ;
  \item une vraie integration de BPM, IAM, GED et IA ;
  \item un effort de demonstrabilite remarquable ;
  \item une progression recente vers plus de clarte metier.
\end{itemize}

\section{Limites actuelles}

Le projet reste toutefois exigeant a stabiliser completement. Plus le perimetre est large, plus les points de jonction sont nombreux. Les principaux points d'attention concernent :

\begin{itemize}[leftmargin=1.2cm]
  \item la maitrise de tous les parcours role par role ;
  \item la precision des statuts metier et de leurs transitions ;
  \item la dependance a une IA locale potentiellement lente ;
  \item la necessite d'un plus grand nombre de tests de non-regression automatisees.
\end{itemize}

\section{Ameliorations metier recommandees}

\begin{itemize}[leftmargin=1.2cm]
  \item structurer explicitement le concept de \texttt{waitingOn} pour savoir si un ticket attend le client, l'agent, le manager ou un tiers ;
  \item imposer des motifs obligatoires pour les pauses SLA, prolongations et escalades ;
  \item enrichir la supervision manager avec une vue file d'attente, charge et risque ;
  \item renforcer la base de connaissances issue des resolutions validees ;
  \item rendre encore plus visibles la prochaine action attendue et la raison de l'etat courant.
\end{itemize}

\section{Ameliorations techniques recommandees}

\begin{itemize}[leftmargin=1.2cm]
  \item etendre les tests d'integration role par role ;
  \item renforcer la gestion centralisee des erreurs API ;
  \item instrumenter davantage la supervision des traitements asynchrones ;
  \item ajouter des tests de parcours end-to-end navigateur sur les scenarios critiques ;
  \item formaliser un contrat d'API de capacites de ticket pour les actions disponibles.
\end{itemize}

\section{Vision d'evolution}

SupportFlow peut evoluer de plusieurs facons. Il peut devenir un outil interne de support IT, un demonstrateur commercial, ou un socle d'application de gestion d'incidents plus large. La presence d'une GED, d'un moteur BPMN et d'une IA locale ouvre aussi des perspectives de capitalisation, de reporting avance et d'automatisation progressive.


